% Redo of story:
% CMA-ES Works
% RL is viable: PPO + Gentle/Clean
% To generalize we use GC-PPO

% Interesting Mention:
% The CMA-Es is more accurate than PPO?? (not always though)
% But optimization combination helps when using PPO for initialization

% Ablations
% About the surrgate and integration with the Latent Decoders
% 


\section{Experiments}
\label{sec:experiments}

\subsection{Is RL Viable for Inverse Fracture Physics?}
\label{sec:rl_viable}

Before developing the goal-conditioned pipeline, we establish that RL can match the standard evolutionary baseline for inverse food fracture. We define two reward functions over the behavioral features that represent different peeling scenarios: \emph{firm peel}, simulating a firm-skinned orange where the objective is clean skin separation with minimal interior disturbance, and \emph{ripe peel}, simulating a ripe orange where the priority is preserving the delicate interior while still achieving adequate separation. These reward functions are weighted sums of percentile-normalized feature terms. CMA-ES is configured with population size $20$, $\sigma_0{=}0.3$ in $[0,1]^9$, and a budget of $5{,}000$ evaluations (sufficient for convergence in this $9$D space). Each per-target PPO agent observes $\mathbf{o}_t = [\mathbf{s}_t,\, \hat{\mathbf{y}}_t,\, R(\hat{\mathbf{y}}_t),\, t/T]$ where $R$ is the fixed reward function (firm or ripe peel), and refines the state over $T{=}8$ steps.

Table~\ref{tab:per_target} reports simulator-validated results. PPO matches or exceeds CMA-ES: PPO Raw achieves $0.712$ on firm peel versus $0.694$ for CMA-ES Raw. CMA-ES VAE collapses to $0.375$ on firm peel, an early indication that unconstrained search in the VAE latent space is vulnerable to surrogate exploitation. This confirms that RL is a viable optimizer for inverse fracture physics, but like CMA-ES, a separate PPO policy must be trained for each new reward function (${\sim}10$ minutes per policy). This limitation motivates the goal-conditioned approach.

\begin{table}[t]
\centering
\caption{Per-target optimization (simulator-validated). Firm and ripe peel represent different peeling scenarios. PPO matches CMA-ES, validating RL for inverse fracture physics. Best in \textbf{bold}, second best \underline{underlined}.}
\label{tab:per_target}
\small
\begin{tabular}{llcc}
\toprule
\textbf{Method} & \textbf{Space} & \textbf{Firm} $\uparrow$ & \textbf{Ripe} $\uparrow$ \\
\midrule
CMA-ES & Raw 9D & $\underline{0.694}$ & $0.725$ \\
CMA-ES & Flow 4D & $0.674$ & $0.725$ \\
CMA-ES & VAE 4D & $0.375$ & $\mathbf{0.731}$ \\
\midrule
PPO & Raw 9D & $\mathbf{0.712}$ & $0.721$ \\
PPO & Flow 4D & $0.674$ & $0.716$ \\
PPO & VAE 4D & $0.682$ & $\underline{0.729}$ \\
\bottomrule
\end{tabular}
\end{table}

\subsection{Goal-Conditioned Inverse Recovery}
\label{sec:gc_exp}

Table~\ref{tab:gc} reports the main result. The top block shows CMA-ES baselines, which optimize independently per target with $5{,}000$ evaluations each. CMA-ES Raw achieves the highest absolute recovery ($0.781$) but at significant per-target cost. CMA-ES in latent spaces suffers from severe surrogate exploitation: CMA-ES VAE achieves only $0.028$ actual recovery despite a surrogate score of $0.498$.

The bottom block shows goal-conditioned policies that generalize across targets through the decode-then-evaluate pipeline (Sec.~\ref{sec:latent}). GC-PPO Flow achieves the best goal-conditioned recovery ($0.642$, $82\%$ of CMA-ES Raw) with $625\times$ fewer evaluations and the smallest gap ($\Delta{=}0.012$) of any method. It outperforms the original $9$D space ($0.523$) by $23\%$, indicating that the flow's learned $4$D manifold provides better structure for RL navigation, consistent with the value discrepancy analysis of Lee~et~al.~\cite{lee2025nfbo}.

All goal-conditioned methods produce estimates in $8$ evaluations (${\sim}10$\,ms) with no per-target retraining. Training takes ${\sim}9$ minutes (one-time) and generalizes to arbitrary targets. By contrast, CMA-ES requires $5{,}000$ evaluations per target, and per-target PPO requires ${\sim}10$ minutes of training for each new reward function. The goal-conditioned policy's training cost is amortized across all future targets, with a break-even point of $3$--$5$ targets compared to per-target CMA-ES.

\begin{table}[t]
\centering
\caption{Inverse recovery on $10$ held-out targets (simulator-validated). Top: CMA-ES baselines (per-target, $5{,}000$ evals each). Bottom: goal-conditioned policies (generalized, $8$ evals each). All GC-PPO variants use the MLP surrogate via decode-then-evaluate.}
\label{tab:gc}
\small
\begin{tabular}{lcccc}
\toprule
\textbf{Method} & \textbf{Evals} & \textbf{Surr.} & \textbf{Actual} $\uparrow$ & $\boldsymbol{\Delta}$ $\downarrow$ \\
\midrule
CMA-ES Raw & $5{,}000$ & $0.922$ & $\mathbf{0.781}$ & $0.141$ \\
CMA-ES Flow & $5{,}000$ & $0.878$ & $0.552$ & $0.327$ \\
CMA-ES VAE & $5{,}000$ & $0.498$ & $0.028$ & $0.470$ \\
\midrule
GC-PPO Raw & $8$ & $0.531$ & $0.523$ & $\mathbf{0.008}$ \\
GC-PPO Flow & $8$ & $0.654$ & $\underline{0.642}$ & $\underline{0.012}$ \\
GC-PPO VAE & $8$ & $0.510$ & $0.485$ & $0.025$ \\
\bottomrule
\end{tabular}
\end{table}

\subsection{Warm-Start Refinement}
\label{sec:warmstart_exp}

The goal-conditioned policy provides instant estimates but is limited by its $8$-step horizon. We refine its output by initializing CMA-ES at the decoded material parameters with $500$ additional evaluations in $[0,1]^9$. Table~\ref{tab:warmstart} reports simulator-validated results alongside two CMA-ES baselines: one with the same $540$-evaluation budget (from random center initialization) and the standard $5{,}000$-evaluation configuration from Table~\ref{tab:gc}. The best warm-start variant (WS VAE, $0.828$) outperforms both CMA-ES baselines. WS Raw ($0.808$) also outperforms both, while WS Flow ($0.773$) underperforms CMA-ES $540$ ($0.788$), showing that the benefit depends on the policy representation. The warm-start gaps ($\Delta{=}0.034$--$0.095$) are smaller than CMA-ES $5{,}000$ ($\Delta{=}0.141$), indicating that policy initialization steers CMA-ES away from surrogate-overestimated regions. All warm-start computations take ${\sim}0.16$\,s on the surrogate, making the per-target cost dominated by the single simulator validation (${\sim}156$\,s).

\begin{table}[t]
\centering
\caption{Warm-start and CMA-ES comparison ($10$ held-out targets, simulator-validated). WS variants use the respective GC-PPO policy for initialization, then CMA-ES refines in $[0,1]^9$.}
\label{tab:warmstart}
\small
\begin{tabular}{lcccc}
\toprule
\textbf{Method} & \textbf{Evals} & \textbf{Surr.} & \textbf{Actual} $\uparrow$ & $\boldsymbol{\Delta}$ $\downarrow$ \\
\midrule
WS Raw & $540$ & $0.874$ & $\underline{0.808}$ & $\underline{0.065}$ \\
WS Flow & $540$ & $0.867$ & $0.773$ & $0.095$ \\
WS VAE & $540$ & $0.862$ & $\mathbf{0.828}$ & $\mathbf{0.034}$ \\
\midrule
CMA-ES Raw & $540$ & $0.867$ & $0.788$ & $0.079$ \\
CMA-ES Raw & $5{,}000$ & $0.922$ & $0.781$ & $0.141$ \\
\bottomrule
\end{tabular}
\end{table}

\subsection{Ablation: Surrogate Configuration}
\label{sec:ablation}

Our default pipeline evaluates all methods through the MLP surrogate ($R^2{=}0.937$) by decoding latent codes back to parameters. An alternative, which is common practice in latent space optimization~\cite{gomez2018automatic, maus2022lolbo, lee2025nfbo}, is to let each latent model use its own prediction head directly, avoiding the decode step. Table~\ref{tab:ablation} shows this degrades performance substantially: GC-PPO Flow drops from $0.642$ to $0.478$ ($-26\%$), and GC-PPO VAE from $0.485$ to $0.458$. The prediction heads have lower accuracy ($R^2{=}0.913$ for the flow head, $0.854$ for the VAE head) because they are trained from compressed $4$D inputs with a multi-task loss rather than from the full $9$D input with a dedicated objective. The gap also increases (flow: $0.012 \to 0.050$), indicating that the lower-accuracy predictor enables more surrogate exploitation. This finding has a practical implication: when latent-space methods are compared against methods using a separate, higher-accuracy evaluator, observed performance differences may reflect evaluator quality rather than representation quality. The shared surrogate approach avoids this by ensuring all methods operate against the same evaluator.

\begin{table}[t]
\centering
\caption{Ablation: surrogate configuration ($10$ held-out targets). Shared: all methods use the MLP surrogate. Own head: each model uses its prediction head.}
\label{tab:ablation}
\small
\begin{tabular}{lccc}
\toprule
\textbf{Method} & \textbf{Surrogate} & \textbf{Actual} $\uparrow$ & $\boldsymbol{\Delta}$ $\downarrow$ \\
\midrule
GC-PPO Flow & Shared MLP & $\mathbf{0.642}$ & $\mathbf{0.012}$ \\
GC-PPO Flow & Own head & $0.478$ & $0.050$ \\
\midrule
GC-PPO VAE & Shared MLP & $\underline{0.485}$ & $0.025$ \\
GC-PPO VAE & Own head & $0.458$ & $\underline{0.023}$ \\
\bottomrule
\end{tabular}
\end{table}